\section{Statements from earlier work}\label[appendix]{app:statements}
This appendix records the supporting statements used in the counter.
Their origins are credited at each statement; proofs are reproduced
in \cref{app:earlier}. Classical normalization, concentration, and
lifting are cited directly to their published sources.

\subsection{Matrix concentration}

\begin{imported}[{Tropp~\cite[Theorem~1.6]{Tropp}}]\label{imp:tropp}
Let \(X_1,\dots,X_Q\) be independent random real \(d_1\times d_2\)
matrices with \(\E X_j=0\) and \(\|X_j\|\le L\) almost surely, and
let
\(\sigma^2=\max\{\|\sum_j\E X_jX_j^{\mathsf T}\|,\|\sum_j\E X_j^{\mathsf T}X_j\|\}\).
Then for all \(x\ge0\),
\[
 \Pr\Bigl[\Bigl\|\sum_jX_j\Bigr\|\ge x\Bigr]
 \le(d_1+d_2)\exp\Bigl(-\frac{x^2/2}{\sigma^2+Lx/3}\Bigr).
\]
\end{imported}

Here and below \(\|\cdot\|\) is the operator norm.

\subsection{Product-star geometry and exact printing}\label{sec:imports-star}

\begin{imported}[{\cite[Lemma~A.1]{Nag72}}]\label{imp:partition}
Let \(2\le r\le n/4\) be even, \(s=\lfloor\lgx(n/r+1)\rfloor\),
\(b=2^s-1\) and \(k=rs\). There are \(2^k\) explicitly described
partitions of \(\{0,1\}^n\), each into \(2^{n-k}\) cells of size
\(2^k\). A point of a cell is specified by one choice from
\(\{0,e_1,\dots,e_b\}\) in each of \(r\) disjoint blocks of
coordinates. For every partition, a uniform cell and a uniform point
give a uniform input. Averaging over the partition and its uniformly
chosen center gives a uniform center, independently of the point
choices. The center bits can be divided into two halves differing in
size by at most one while keeping each block on its assigned side.
\end{imported}

Since \(2^{s+1}>n/r+1\), we have \(b+1>(n/r+1)/2>n/(2r)\). We write a
cell's address, or \emph{column}, as \(z\in\{0,1\}^{n-k}\) and its
center as \(c(z)\). The \(r/2\) blocks on the left side give \(N=2^{k/2}\)
left patterns \(u\in[N]\), and the \(r/2\) blocks on the right side
give \(N\) right patterns \(v\in[N]\); we identify a pattern with its
error vector, so the points of the cell are \(c(z)\oplus u\oplus v\).
Let \(F(u)\) be the set of at most \(r/2\) coordinates flipped by
\(u\); it depends only on \(u\) and the partition, not on \(z\).

\begin{imported}[{\cite[Lemma~A.2]{Nag72}}]\label{imp:range}
For every threshold function \(T\) and all \(c,u,v\) as above,
\(\Delta_T(c,u,v)=T(c\oplus u\oplus v)-T(c\oplus u)-T(c\oplus v)+T(c)\in\{-1,0,1\}\).
\end{imported}

\begin{imported}[{\cite[Theorem~A.4]{Nag72}}]\label{imp:population}
Average over a uniform partition, a uniform center, and uniform
patterns \(u,v\). For a threshold function \(T\) that depends on \(f\)
coordinates, \(\Pr[\Delta_T\ne0]\le2f/(b+1)^{3/2}\).
\end{imported}

The next four statements compile threshold atoms into tables
printable over all columns.

\begin{imported}[{\cite[Lemma~C.2]{Nag72}}]\label{imp:disjoint}
Let \(H(x)=\ind[\sum_{i=1}^fw_ix_i>T]\) have integer parameters, let
\(B_0'=\max(1,|T|,|w_1|,\dots,|w_f|)\) and
\(\ell=\max\{1,\lceil\lgx((f+1)B_0'+1)\rceil\}\). There is an explicitly
constructible list of at most \(\max(1,f\ell)\) integer affine
equations \(E_j(x)=0\) with \(H(x)=\sum_j\ind[E_j(x)=0]\) on the
Boolean cube, the sum always lying in \(\{0,1\}\). Their coefficients
have \(O(\ell+\log(f+2))\) bits. Weak thresholds are included by
replacing \(\sum_iw_ix_i\ge T\) with \(\sum_iw_ix_i>T-1\).
\end{imported}

\begin{imported}[{\cite[Lemma~C.3 and the paragraph after it]{Nag72}}]\label{imp:pullback}
In the product-star parametrization, with each block kept on its side
of the column split \(z=(z_{\rm L},z_{\rm R})\), each original input
coordinate of \(c(z)\oplus x\), for a fixed point \(x\) of the cell, is
a Boolean function, computable in polynomial time, of one of the two
column halves. Substituting these functions into a disjoint exact
child \(c_j+\sum_id_{j,i}x_i=0\) of \cref{imp:disjoint} yields a
\emph{separated} equation \(a_j(z_{\rm L})+b_j(z_{\rm R})=0\) with
polynomial-bit values; the children remain disjoint and represent the
pulled-back threshold exactly over the integers.
\end{imported}

\begin{imported}[{\cite[Lemma~C.6]{Nag72}}]\label{imp:stack}
If integer separated equations satisfy \(|E_j(z)|\le V\) for all
\(j,z\) and \(H=2V+2\), then
\(\prod_{j=1}^t\ind[E_j(z)=0]=\ind[\sum_jH^{j-1}E_j(z)=0]\), and the
stacked equation is separated with \(O(t\log(V+2))\)-bit
coefficients.
\end{imported}

\begin{imported}[{\cite[Theorem~C.5]{Nag72}}]\label{imp:printer}
Fix a polynomial description envelope and \(n/2\le n_{\rm out}\le n\).
Let \(E_i(z)=a_i(z_{\rm L})+b_i(z_{\rm R})\), \(1\le i\le s\), be
separated equations over \(z\in\{0,1\}^{n_{\rm out}}\) with
polynomial-bit half-scores computable in polynomial time, and let
\(v_i\) be polynomial-bit integers. If \(s\le2^{n_{\rm out}/100}\),
the table \(z\mapsto\sum_{i=1}^sv_i\ind[E_i(z)=0]\) is printable in
\(2^{n_{\rm out}}\poly(n)\) bit operations, in addition to the
construction of the list of the \(s\) equations.
\end{imported}

The printer reduces the table to weighted dominance counting, which
Matou\v{s}ek's bucketing~\cite{Mat91} reduces to Williams' uniform
rectangular matrix multiplication~\cite[Corollary~4.4]{WilACC14}.

\begin{imported}[{\cite[Lemma~F.2 and its proof]{Nag72}}]\label{imp:seed}
Let \(N=2^{k/2}\), let \(Q\) and \(h\) be positive integers, take an
integer \(\ell_f\ge\max\{k,\lceil\lgx(Q+1)\rceil\}\) and the field
\(\F_{2^{\ell_f}}\), label the \(Q\) samples by distinct field elements,
evaluate a uniformly random polynomial of degree \(2h-1\) at the
labels, and project each value to \(k\) bits, read as a pair in
\([N]^2\). The seed has \(S=2h\ell_f\) bits, and the \(Q\) sample
positions are \(2h\)-wise independent and uniform on \([N]^2\). For
every seed prefix, the conditional probability of any conjunction of
events ``sample \(j\) lies at position \(x\)'' is \(0\) or \(2^{-a}\)
with \(a\le S\), and it is computable by Gaussian elimination over
\(\F_2\); repeated sample labels repeat or augment the same linear
system. The two one-bit extensions of a prefix average to the prefix:
for every function \(\Psi\) of the seed,
\(\E[\Psi\mid\pi]=\tfrac12(\E[\Psi\mid\pi0]+\E[\Psi\mid\pi1])\).
\end{imported}

\subsection{Normalization, sampling rates, and residue reconstruction}\label{sec:imports-nc}

\begin{imported}[{\cite[Lemma~3.2]{CTW26}}]\label{imp:integer}
Every threshold function on \(f\ge1\) Boolean inputs has an
equivalent representation whose weights and threshold are integers in
\(\pm f^f\).
\end{imported}

This is~\cite[Lemma~3.2]{CTW26}, attributed there to Muroga, Toda, and
Takasu~\cite{MTT61}.

\begin{imported}[{\cite[Theorem~2.8.4]{Vershynin18}}]\label{imp:bernstein}
If independent centered random variables \(Y_t\) satisfy
\(|Y_t|\le1\) and \(\sum_t\E Y_t^2\le V\), then for every \(T>0\),
\(\Pr[|\sum_tY_t|\ge\sqrt{2TV}+2T/3]\le2e^{-T}\).
\end{imported}

\begin{imported}[{\cite[Lemma~5.1 and Equation~(1)]{NagCG}}]\label{imp:chain}
Let \(T>0\), let \(h\ge64T\) be an integer, and let \(\ell_0\ge0\).
Define
\[
 \bar\ell=\max\{\ell_0,h^2/(256T)\},\quad
 \pi^*=\frac{h^2}{8T\bar\ell},\quad
 s(\pi)=\sqrt{2T\pi\bar\ell}+\frac{2T}3,\quad
 M_1=\lceil4\bar\ell/h\rceil,\quad\pi_1=1/M_1 .
\]
While \(x_r=\pi_r/\pi^*<2^{-8}\), set \(a_r=\lfloor1/(4\sqrt{x_r})\rfloor\)
and \(\pi_{r+1}=a_r\pi_r\); stop at the first \(J\) with
\(x_J\ge2^{-8}\). Then:
\begin{enumerate}[(i)]
\item \(a_r\ge4\), all rates lie in \((0,1]\), and \(\pi_r=\prod_{i<r}a_i/M_1\);
\item \(\pi_1\bar\ell+s(\pi_1)\le h/2\), and
\(s(\pi_{r+1})+a_rs(\pi_r)\le h/2\) for \(r<J\);
\item \(J\le\lceil\lgx\lgx(h/T)\rceil+1\);
\item \(s(\pi_J)/\pi_J\le86T\bar\ell/h\).
\end{enumerate}
Put \(A=\prod_{i=1}^{J-1}a_i\), \(A_0=A\) and
\(A_r=\prod_{i=r+1}^{J-1}a_i\) for \(1\le r<J\). For arbitrary integer
vectors \(K_1,\dots,K_J\), with \(E_0=K_1\) and
\(E_r=K_{r+1}-a_rK_r\), the identities
\(K_J=\sum_{r=0}^{J-1}A_rE_r\) and \(\pi_J=A/M_1\) hold.
\end{imported}

The telescoping identity holds for \emph{any} proposed vectors, not
only for nested samples; \cref{lem:certify}(d) uses it for arbitrary witness vectors.

\begin{imported}[{\cite[Lemmas~4.1--4.4 and the paragraph after Lemma~4.3]{NagCG}}]\label{imp:window}
Let \(\cP\) be the set of primes in \((P/2,P]\) and
\(D'=1+\sum_{p\in\cP}(p-1)\). For a finitely supported
\(F:\Z^d\to\Q\) and a prime \(p\), let
\(F_p(\rho)=\sum_{y\equiv\rho\ (\mathrm{mod}\ p)}F(y)\) for
\(\rho\in(\Z/p)^d\).
\begin{enumerate}[(a)]
\item A rational function supported on at most \(D'\) consecutive
integers with \(F_p=0\) for all \(p\in\cP\) is zero.
\item Let \(S_\varphi(q)=\sum_{t\le q}t\varphi(t)\), with \(\varphi\)
Euler's function, and let \(L_2\) be the largest \(L\ge2\) with
\(S_\varphi(\lfloor(D'-1)/(L-1)\rfloor)\ge L\), or \(L_2=1\) if there is
none. A rational function supported on a given box in \(\Z^2\) with
both sides at most \(L_2\) is uniquely determined, among such
functions, by its push-forwards \(F_p\), \(p\in\cP\).
\item As \(P\to\infty\), \(D'=\Theta(P^2/\ln P)\) and
\(L_2=\Theta(P^{3/2}/(\ln P)^{3/4})\).
\end{enumerate}
Exact Gaussian elimination, with one unknown per box point and one
equation per residue class, finds the unique solution when the data
come from a function on the box.
\end{imported}

Part~(b) is the quantitative modular form, proved
in~\cite{NagCG}, of Hajdu and Tijdeman's factorization method for
discrete tomography~\cite[Theorem~1]{HT01}. If \(L>D'\) then
\(\lfloor(D'-1)/(L-1)\rfloor=0\) and \(S_\varphi(0)=0<L\), so
\(L_2\le D'\).

\begin{imported}[{\cite[Section~2.3]{BT94}}]\label{imp:amp}
For \(e\ge1\),
\[\Amp_e(t)=1-(1-t)^e\sum_{j=0}^{e-1}\binom{e-1+j}jt^j\]
is an integer polynomial of degree \(2e-1\). Modulo \(p^e\), it maps
every integer congruent to \(0\) modulo \(p\) to \(0\) and every integer
congruent to \(1\) modulo \(p\) to \(1\).
\end{imported}

This is the modulus-amplifying polynomial of Beigel and
Tarui~\cite[Section~2.3]{BT94}.

